% Part: proof-theory
% Chapter: proof-search
% Section: introduction

\documentclass[../../../include/open-logic-section]{subfiles}

\begin{document}

\olfileid[hi]{pt}{ps}{int}

\olsection{परिचय}

\Log{G3c} या \Log{N1i} जैसी !!{proof} प्रणालियों के अभिगृहीत और अनुमान-नियम
यह निर्धारित करते हैं कि अनुक्रमों अथवा !!{formula}s के कौन-से वृक्ष सही
!!{proof}s माने जाएँगे। अनुक्रमों या !!{formula}s का ऐसा कोई वृक्ष दिए जाने
पर (संभवतः इस अतिरिक्त सूचना सहित कि प्रत्येक चरण में कौन-से नियम लागू हुए
हैं, अथवा मान्यताओं और उन्हें मुक्त करने वाले अनुमान-नियमों के चिह्नों सहित),
अभिगृहीतों और नियमों की परिभाषाएँ हमें जाँचने देती हैं कि दिया गया वृक्ष सही
!!a{proof} है या नहीं। दी हुई मान्यताओं से किसी दिए हुए अनुक्रम या दिए हुए
!!{formula} का !!a{proof} \emph{खोजना} एक अलग और प्रायः अधिक कठिन समस्या है।
यही \emph{प्रमाण-खोज} की समस्या है।

प्रमाण-खोज की एक अत्यंत सीधी-सादी प्रक्रिया है। हम !!{formula}s को किसी सीमित
वर्णमाला के प्रतीकों के अनुक्रम मान सकते हैं। अनुक्रमों और अनुक्रमों अथवा
!!{formula}s के वृक्षों को भी !!{formula}s के अनुक्रम माना जा सकता है (वृक्ष
की शाखाओं को दिखाने के लिए शायद विशेष कोष्ठकों का उपयोग करके)। अभिगृहीत और
नियम हमें यांत्रिक ढंग से जाँचने देते हैं कि प्रतीकों का कोई अनुक्रम सही
!!a{proof} निरूपित करता है या नहीं। फलतः हम \emph{सभी संभव} सही !!{proof}s
यांत्रिक ढंग से उत्पन्न कर सकते हैं (!!{proof}s को निरूपित करने वाली वर्णमाला
के प्रतीकों के सभी अनुक्रम बनाकर और प्रत्येक प्रत्याशी अनुक्रम को जाँचकर)।
सीधी-सादी प्रमाण-खोज प्रक्रिया यह है: सभी सही !!{proof}s तब तक उत्पन्न करें
जब तक दिए हुए अनुक्रम का (या दिए हुए !!{formula} का ऐसा प्रमाण जिसकी खुली
मान्यताएँ दी हुई मान्यताओं में हों) कोई !!a{proof} न मिल जाए।

इससे अच्छी विधि वही सहज विधि है जिसका प्रयोग आप हाथ से !!a{proof} खोजने में
करेंगे: अंतिम अनुक्रम से आरंभ करें, कोई !!a{formula} चुनें और किसी अनुमान-नियम
को ``उलटी दिशा में'' लागू करके ऐसा आधार-वाक्य या ऐसे आधार-वाक्य उत्पन्न करें
जिनके !!{proof}s मिलने पर वे अंतिम अनुमान के साथ मिलकर दिए हुए अनुक्रम का
!!a{proof} बना दें। प्राकृतिक निगमन प्रणाली में !!a{proof} खोजने के लिए भी
इसी प्रकार की (यद्यपि अधिक जटिल) विधि अपनाई जाती है।

पर यह विधि अचूक नहीं है: गलत नियम या !!{formula}s का गलत क्रम चुनने पर आप
ऐसी स्थिति में फँस सकते हैं जहाँ आगे कोई मार्ग न हो। यह विधि पूरी तरह
विनिर्दिष्ट भी नहीं है: प्रत्येक बिंदु पर चुनने योग्य एक से अधिक !!{formula}s
या उलटी दिशा में लागू करने योग्य एक से अधिक नियम हो सकते हैं। \emph{प्रमाण-खोज
कलन-विधि} पूरी तरह विनिर्दिष्ट ऐसी प्रक्रिया है जो !!a{proof} विद्यमान होने पर
उसे खोज लेगी (अर्थात् जो अचूक है)।

कुछ !!{proof} प्रणालियाँ अन्य प्रणालियों की अपेक्षा सरल प्रमाण-खोज
कलन-विधियों के अधिक अनुकूल होती हैं। अनुक्रम प्रणालियों में किसी !!{formula}
का !!{main operator} और उसका बाईं या दाईं ओर आना यह निर्धारित करता है कि
उलटी दिशा में कौन-सा नियम लगाया जाए। उदाहरणतः यदि आप $\Gamma \Sequent
\Delta, !A \land !B$ का !!a{proof} खोजना चाहते हैं, जहाँ $!A \land !B$ दाईं
ओर है, तो \RightR{\land} को उलटी दिशा में लगाकर दो संगत आधार-वाक्य $\Gamma
\Sequent \Delta, !A$ और $\Gamma \Sequent \Delta, !B$ उत्पन्न करने चाहिए।
लेकिन यदि आपकी प्रणाली \Log{G1c} है, तो संकुचन की उपस्थिति अधिक विकल्प देकर
काम कठिन कर देती है: आप \RightR{\Contraction} को भी उलटी दिशा में लगाकर
$\Gamma \Sequent \Delta, !A \land !B, !A \land !B$ आधार-वाक्य बना सकते हैं;
फिर $\Gamma \Sequent \Delta, !A \land !B, !A \land !B, !A \land !B$, और
इसी प्रकार आगे। \Log{G2c} में स्थिति और भी खराब है। \RightR{\land} को उलटी
दिशा में लगाने के लिए आपको ऐसे $\Gamma_1$, $\Gamma_2$ का अनुमान लगाना होगा कि
$\Gamma = \Gamma_1 \cup \Gamma_2$, और ऐसे $\Delta_1$, $\Delta_2$ का कि
$\Delta = \Delta_1 \cup \Delta_2$, फिर $\Gamma_1 \Sequent \Delta_1, !A$
और $\Gamma_2 \Sequent \Delta_2, !B$ आधार-वाक्यों पर भाग्य आजमाना होगा। गलत
अनुमान आगे जाकर मार्ग बंद कर सकता है, तब आरंभ तक लौटकर दूसरे $\Gamma_1$,
$\Gamma_2$, $\Delta_1$, $\Delta_2$ चुनने पड़ेंगे। यदि !!{formula} $!A \lor
!B$ है, तो \RightR{\lor} की दो संभावनाओं में से एक चुनकर $\Gamma \Sequent
\Delta, !A$ या $\Gamma \Sequent \Delta, !B$ में से एक आधार-वाक्य उत्पन्न
करना होगा। गलत चुनाव पर फिर पीछे लौटना पड़ेगा।

प्रणाली \Log{G3c} में ये समस्याएँ नहीं हैं। उसमें कोई संरचनात्मक नियम नहीं
है, और अनुमान-नियम का चुनाव !!{main operator} तथा !!{formula} के आने वाले पक्ष
से निर्धारित होता है। अतः \Log{G3c}, \Log{G1c} या \Log{G2c} की अपेक्षा
प्रमाण-खोज के लिए अधिक उपयुक्त है। सफल खोज \Log{G3c} में !!a{proof} देगी और
उस !!{proof} का फिर \Log{G1c} या \Log{G2c} (या इस दृष्टि से \Log{N1c}) में
अनुवाद किया जा सकता है। इसलिए \Log{G3c} की प्रमाण-खोज कलन-विधि और \Log{G3c}
के !!{proof}s को अन्य प्रणालियों के !!{proof}s में बदलने वाली अनुवाद
कलन-विधियाँ मिलकर उन प्रणालियों की प्रमाण-खोज समस्या भी हल करती हैं।

हम कैसे जानें कि कोई प्रमाण-खोज कलन-विधि अचूक है? इसके लिए दिखाना होगा कि यदि
कलन-विधि !!a{proof} खोजने में विफल रहती है, तो कोई !!a{proof} हो ही नहीं सकता।
आखिर कोई खराब कलन-विधि चुनने पर ऐसे प्रसंग आ सकते हैं जहाँ प्रमाण होते हुए भी
वह उसे न खोज पाए। सीधी-सादी कलन-विधि में यह समस्या नहीं है, क्योंकि वह सभी
संभव !!{proof}s में खोज करती है। किंतु प्रमाण-खोज कलन-विधि से अपेक्षा है कि वह
दक्ष हो और न्यूनतम प्रयास करे।

प्रमाण-खोज कलन-विधि को अचूक दिखाने की मुख्य युक्ति यह है कि उसकी विफलता के
ढंग से ही सिद्ध कर दें कि उस अनुक्रम का कोई !!a{proof} संभव नहीं है। इसके लिए
दिखाते हैं कि अनुक्रम वैध नहीं है। चिरसम्मत प्रथम-कोटि तर्क में वैध अनुक्रम
$\Gamma \Sequent \Delta$ वे हैं जिनमें प्रत्येक !!{structure}, $\Gamma$ के
कम-से-कम एक !!{formula} को असत्य या~$\Delta$ के कम-से-कम एक !!{formula} को
सत्य बनाती है। \Log{G3c} \emph{निर्दोष} है, अर्थात् वह केवल वैध अनुक्रमों को
!!{prove} करती है। यह !!{proof}s पर आगमन से स्थापित होता है: अभिगृहीत स्पष्टतः
वैध हैं, और किसी नियम के आधार-वाक्य वैध हों तो उसका निष्कर्ष भी वैध है। किसी
प्रमाण-खोज कलन-विधि को अचूक दिखाने के लिए हम सिद्ध करते हैं कि यदि $\Gamma
\Sequent \Delta$ के !!a{proof} की खोज विफल होती है, तो ऐसी !!a{structure}
परिभाषित की जा सकती है जिसमें $\Gamma$ के सभी !!{formula}s सत्य और~$\Delta$
के सभी !!{formula}s असत्य हों। इससे \Log{G3c} की \emph{पूर्णता} भी स्थापित
होती है: यदि $\Gamma \Sequent \Delta$ वैध है, तो उसका कोई
\Log{G3c}-!!{proof} है। प्रत्येक विफल प्रमाण-खोज $\Gamma \Sequent \Delta$ को
अवैध दिखाती है, इसलिए वैध अनुक्रम की प्रत्येक प्रमाण-खोज सफल होनी चाहिए।

\end{document}
